\documentclass[11pt]{article}
\usepackage[T1]{fontenc}
\usepackage{lmodern,microtype}
\usepackage[margin=1in]{geometry}
\usepackage{amsmath,amssymb,amsthm,mathtools}
\usepackage{booktabs,longtable,array,xcolor,xurl}
\usepackage[colorlinks=true,linkcolor=blue!40!black,citecolor=blue!40!black,urlcolor=blue!45!black]{hyperref}
\usepackage{bookmark,fancyhdr}
\setlength{\headheight}{14pt}
\pagestyle{fancy}\fancyhf{}
\fancyhead[L]{\small Stieltjes correlations and spectral jets}
\fancyhead[R]{\small ProveIt research continuation}
\fancyfoot[C]{\thepage}
\setlength{\parskip}{3pt}
\emergencystretch=2em
\allowdisplaybreaks[2]
\numberwithin{equation}{section}
\newtheorem{theorem}{Theorem}[section]
\newtheorem{proposition}[theorem]{Proposition}
\newtheorem{lemma}[theorem]{Lemma}
\newtheorem{corollary}[theorem]{Corollary}
\theoremstyle{definition}
\newtheorem{definition}[theorem]{Definition}
\newtheorem{example}[theorem]{Example}
\theoremstyle{remark}
\newtheorem{remark}[theorem]{Remark}
\DeclareMathOperator{\Li}{Li}
\DeclareMathOperator{\FP}{FP}
\DeclareMathOperator{\Rea}{Re}
\DeclareMathOperator{\Ima}{Im}
\newcommand{\dd}{\,\mathrm d}
\newcommand{\C}{\mathbb C}
\newcommand{\Q}{\mathbb Q}
\newcommand{\Zalg}{\mathcal Z}
\newcommand{\gzero}{\gamma}
\newcommand{\coeff}[1]{\left[#1\right]}
\newcommand{\LiJ}[2]{\left.\partial_s^{#1}\Li_s(#2)\right|_{s=2}}
\newcommand{\file}[1]{\nolinkurl{#1}}
\title{Bilinear Closure of Generalized\ Stieltjes Correlations\\[5pt]
\large Gamma-ratio kernels, polylogarithmic jets,\ exact collision expansions and antiderivatives}
\author{Research continuation prepared for Vladimir Reshetnikov's ProveIt project\\
\small Mathematical derivations and computational audit with OpenAI assistance}
\date{10 October 2026}
\hypersetup{pdftitle={Bilinear Closure of Generalized Stieltjes Correlations},pdfauthor={ProveIt research continuation with OpenAI assistance}}
\begin{document}
\maketitle
\thispagestyle{plain}
\begin{abstract}
We study products of generalized Stieltjes functions under a circular
translation. Their ordinary integrals diverge, but a specified endpoint
finite part admits a holomorphic two-parameter generating function.
A Gamma-ratio identity proves that every bilinear finite-part correlation
is a finite \emph{linear} combination of single Stieltjes values at the
translation and its reflection, with coefficients polynomial in positive
integer zeta values. The highest Stieltjes order is explicit, and the next
order is absent identically. We derive the coincident-endpoint analogue,
an exact convergent collision expansion with a canonical logarithmic
counterterm, rational-grid and character identities, and explicit
antiderivatives. An ordinary, nonregularized log-Gamma autocorrelation
also reduces to Hurwitz-zeta derivatives at $-1$, equivalently to spectral
polylogarithm derivatives at $2$.

The proofs use classical Hurwitz Fourier theory and local subtraction;
the finite-part closure and its consequences are developed here as a
continuation, not asserted to have established literature priority.
The package contains exact coefficient identities through total order
$8$, $272$ symbolic checks, and $22$ independent numerical diagnostics.
A reproducible numerical differentiation canary is recorded separately.
Neither floating-point agreement nor the symbolic checks replace the
analytic proofs, and no proof-assistant formalization is claimed.
\end{abstract}
\newpage
\tableofcontents
\newpage

\section{Research position and the main result}
\label{sec:position}

The ProveIt manuscript connects polylogarithms, Gamma values, Hurwitz-zeta
jets and generalized Stieltjes functions. Its integration chapter already
contains the classical single-function ladder that expresses an
antiderivative of $\gamma_n(a)$ through $\zeta^{(n+1)}(0,a)$
\cite{repo-integration}. A natural next question is whether a product of two
Stieltjes functions introduces genuinely bilinear special-function
coordinates after integration. For the translated finite parts studied
here, the answer is particularly simple: \emph{it does not}.

Here and below, a superscript on $\zeta$ differentiates its spectral
variable. The functions $\gamma_n(a)$ are fixed by
\begin{equation}
 \zeta(1+u,a)=\frac1u+
 \sum_{n\geq0}\frac{(-1)^n}{n!}\gamma_n(a)u^n,
 \qquad \gamma_0(a)=-\psi(a).
 \label{eq:stieltjes-definition}
\end{equation}
Write $b=1-a$ for $0<a<1$, and let $I_{mn}(a)$ denote the endpoint finite
part specified in Definition~\ref{def:finitepart}. Our principal result is
\begin{equation}
\begin{split}
 I_{mn}(a)={}&\frac{\gamma_{m+n+1}(a)}{m+1}
       +\frac{\gamma_{m+n+1}(1-a)}{n+1}+e_{mn}\\
 &+\sum_{j=0}^{m+n-1}
       \bigl(c_{mn,j}\gamma_j(a)+d_{mn,j}\gamma_j(1-a)\bigr),
 \label{eq:front-result}
\end{split}
\end{equation}
where the sum is empty if $m+n=0$, and all coefficients belong to
$\Q[\zeta(2),\ldots,\zeta(m+n+2)]$. In particular there is no term of
order $\gamma_{m+n}$. The statement is an identity for every $m,n\geq0$,
not an extrapolation from the accompanying finite table.

At the first order, the result reads
\begin{equation}
 \FP\int_0^1\psi(x)\psi(x\mathbin\oplus a)\dd x
 =\gamma_1(a)+\gamma_1(1-a)-\frac{\pi^2}{3}.
 \label{eq:first-result}
\end{equation}
Here $x\mathbin\oplus a$ is addition modulo one, interpreted on the open
pieces of the interval. The symbol $\FP$ is essential. Removing it makes
\eqref{eq:first-result} false as a statement about an ordinary integral.
An equivalent absolutely convergent form is given in
\eqref{eq:digamma-convergent}.

The relevant classical background is the Hurwitz Fourier expansion,
its parameter-shift law, and the Lerch identity
\cite{dlmf-hurwitz,dlmf-polylog}. Integrals of products of Hurwitz zeta
functions have substantial prior literature; in particular
Shpot--Chaudhary--Paris treat products of the auxiliary functions
$\zeta(s,1+x)$ and related reflected products \cite{shpot}.
We do not claim the Fourier correlation mechanism itself is new.
The specific all-order finite-part closure, its missing-order rule,
and the collision-compatible coefficient presentation are the research
outputs proved in this article. A comprehensive priority determination
would require a wider literature comparison than this delivery provides.

\subsection{What is, and is not, settled}
This continuation proves exact identities rather than addressing the
recent zero-geometry and extremizer problems. It does not prove numerical
independence of the resulting constants. Nor does it settle the remaining
colored-polylogarithm reduction problems in the manuscript.
The inspected current $S_4$ chapter already contains an exact proof and
certificate \cite{repo-s4}; it must not be reclassified as an open
conjecture on the basis of older research notes. The surviving $S_6$
reduction problem recorded in the discovery chapter is not resolved here
\cite{repo-discovery}.

The repository review was targeted: selected manuscript sections, source
metadata, and the incoming directory inventory and intake instructions
were inspected. The incoming binary ZIP contents were not audited.
No existing repository theorem is declared false in this report.
Section~\ref{sec:audit} identifies normalization and computational hazards
that an integration of these new formulas must avoid.

\section{Analytic conventions and a translated Hurwitz kernel}
\label{sec:kernel}

For $0<a<1$, set $b=1-a$ and define
\begin{equation}
\begin{split}
 C(s,t;a)={}&\int_0^b\zeta(s,x)\zeta(t,a+x)\dd x\\
 &+\int_0^a\zeta(s,b+y)\zeta(t,y)\dd y.
 \label{eq:C-definition}
\end{split}
\end{equation}
This is the circularly translated product, without complex conjugation.
All powers of positive real quantities use their real logarithms.
At $e^{\pm2\pi ia}$ we use the standard polylogarithm continuation in its
argument, with spectral derivatives taken on that fixed branch.
Since $0<a<1$, these arguments are never $1$.

\begin{lemma}[Convergence and zero mean]
\label{lem:convergence}
The integral \eqref{eq:C-definition} is absolutely convergent and jointly
holomorphic for $\Rea s<1$, $\Rea t<1$. In this region,
\begin{equation}
 \int_0^1\zeta(s,x)\dd x=0.
 \label{eq:zero-mean}
\end{equation}
\end{lemma}
\begin{proof}
Use $\zeta(s,x)=x^{-s}+\zeta(s,1+x)$.
On the first piece only the factor with parameter $x$ is singular at
zero; on the second piece only the factor with parameter $y$ is singular.
On compact subsets of the stated spectral region, the other factors and
all their spectral derivatives are bounded. The potentially singular
terms are dominated by $x^{-\sigma}(1+|\log x|^k)$ with $\sigma<1$,
or the corresponding expression in $y$. This proves convergence,
holomorphy and differentiation under the integral on compact subsets.
Finally, the parameter derivative identity
\begin{equation}
 \partial_x\zeta(s-1,x)=(1-s)\zeta(s,x)
 \label{eq:hurwitz-derivative}
\end{equation}
and the shift law show that the endpoint values of
$\zeta(s-1,x)$ at zero and one agree when $\Rea s<1$.
Integration proves \eqref{eq:zero-mean}.
\end{proof}

\begin{theorem}[Two forms of the translated kernel]
\label{thm:kernel}
For $0<a<1$ and $\Rea s,\Rea t<1$, put $z=e^{2\pi ia}$.
Then
\begin{align}
 C(s,t;a)={}&\Gamma(1-s)\Gamma(1-t)(2\pi)^{s+t-2}
 \bigl\{e^{i\pi(s-t)/2}\Li_{2-s-t}(z^{-1})\notag\\
 &\hspace{42mm}+e^{-i\pi(s-t)/2}\Li_{2-s-t}(z)\bigr\},
 \label{eq:C-polylog}\\
 C(s,t;a)={}&\Gamma(s+t-1)
 \left\{\frac{\Gamma(1-t)}{\Gamma(s)}\zeta(s+t-1,b)
       +\frac{\Gamma(1-s)}{\Gamma(t)}\zeta(s+t-1,a)\right\}.
 \label{eq:C-hurwitz}
\end{align}
All apparent singularities of the complete right side in this region
are interpreted by analytic continuation. The formulas also give a
meromorphic continuation in the two spectral variables.
\end{theorem}
\begin{proof}
For $\Rea s<0$, the classical Hurwitz Fourier series is absolutely and
uniformly convergent, and its nonzero Fourier coefficients are
\begin{equation}
 \widehat\zeta_s(n)=
 \frac{\Gamma(1-s)}{(2\pi)^{1-s}|n|^{1-s}}
 \exp\!\left(-\frac{i\pi(1-s)}2\operatorname{sgn}n\right),
 \qquad n\ne0.
 \label{eq:fourier-coefficients}
\end{equation}
The zero coefficient is zero. This is the exponential form of the
Hurwitz Fourier identity \cite[Eq.~25.11.9]{dlmf-hurwitz}.
Multiplying the two absolutely convergent Fourier series and integrating
leaves only frequencies $n$ and $-n$. Translation contributes the phase
$e^{-2\pi ina}$. The sum over $n>0$ is the first polylogarithm in
\eqref{eq:C-polylog}, and the sum over $n<0$ is the second.

For completeness, write $r=s+t-1$ and insert the Hurwitz Fourier
expansions of $\zeta(r,a)$ and $\zeta(r,b)$ into
\eqref{eq:C-hurwitz}. Gamma reflection reduces the coefficient of
$\Li_{1-r}(z)$ to the desired phase, using
\[
 \frac{\sin(\pi s)e^{-i\pi(s+t)/2}
       +\sin(\pi t)e^{i\pi(s+t)/2}}{\sin\pi(s+t)}
 =e^{-i\pi(s-t)/2}.
\]
The other coefficient follows by reversing the phases. Initially avoid
integer hyperplanes where this manipulation divides by zero. The
result on their complement extends across removable singularities.
The identity theorem and Lemma~\ref{lem:convergence} now extend both
expressions to $\Rea s,\Rea t<1$. Their right sides supply the asserted
meromorphic continuation.
\end{proof}

\begin{remark}[Translation is not ordinary reflection]
The two-piece definition is not
$\int_0^1\zeta(s,x)\zeta(t,1-x)\dd x$ and is not
$\int_0^1\zeta(s,x)\zeta(t,a+x)\dd x$ without wrapping.
Replacing one by another changes the Fourier phase and the endpoint
terms. Also, $a=0$ is not a harmless substitution in the ordinary integral:
the two singularities then coincide. A sufficient convergence region at
$a=0$ includes $\Rea(s+t)<1$ as well as $\Rea s,\Rea t<1$.
Section~\ref{sec:coincident} handles the required continuation explicitly.
\end{remark}

\section{Endpoint subtraction and the Stieltjes generating function}
\label{sec:finitepart}

The Hurwitz shift law gives the exact local decomposition
\begin{equation}
 \gamma_m(x)=\frac{\log^m x}{x}+\gamma_m(1+x),
 \qquad x>0.
 \label{eq:stieltjes-shift}
\end{equation}
This identity determines the counterterms without a numerical fitting
procedure or a choice of divergent Fourier summation.

\begin{definition}[Separated-endpoint finite part]
\label{def:finitepart}
For integers $m,n\ge0$ and $0<a<1$, define
\begin{align}
 I_{mn}(a)=\lim_{\epsilon\downarrow0}\biggl\{
 &\int_\epsilon^b\gamma_m(x)\gamma_n(a+x)\dd x
 +\int_\epsilon^a\gamma_m(b+y)\gamma_n(y)\dd y\notag\\
 &+\frac{\gamma_n(a)}{m+1}\log^{m+1}\epsilon
 +\frac{\gamma_m(b)}{n+1}\log^{n+1}\epsilon\biggr\}.
 \label{eq:I-cutoff}
\end{align}
The cutoff is smaller than $\min(a,b)$. Equivalently the two endpoints
can be assigned separate cutoffs, with their respective counterterms.
\end{definition}

\begin{lemma}[An absolutely convergent implementation]
\label{lem:FP-exists}
The limit exists and equals
\begin{align}
 I_{mn}(a)={}&\int_0^b\left\{\gamma_m(x)\gamma_n(a+x)
               -\gamma_n(a)\frac{\log^m x}{x}\right\}\dd x\notag\\
 &+\int_0^a\left\{\gamma_m(b+y)\gamma_n(y)
               -\gamma_m(b)\frac{\log^n y}{y}\right\}\dd y\notag\\
 &+\frac{\gamma_n(a)}{m+1}\log^{m+1}b
  +\frac{\gamma_m(b)}{n+1}\log^{n+1}a.
 \label{eq:I-convergent}
\end{align}
It also satisfies $I_{mn}(a)=I_{nm}(1-a)$.
\end{lemma}
\begin{proof}
On the first interval, use \eqref{eq:stieltjes-shift}. The singular
remainder is
$\log^m x\,[\gamma_n(a+x)-\gamma_n(a)]/x$, which is integrable because
the bracket is $O(x)$. The other term is bounded there. The argument at
the second endpoint is identical. Integrating the explicitly subtracted
$\log^m x/x$ and $\log^n y/y$ terms gives
\eqref{eq:I-convergent}. Interchange the two intervals for the symmetry.
\end{proof}

Set
\begin{equation}
 G(u,x)=\zeta(1+u,x)-\frac1u,
 \qquad G(0,x)=-\psi(x),\qquad g(u)=\Gamma(1-u)(2\pi)^u.
 \label{eq:G}
\end{equation}
The removable value of $G$ at zero is always included. Define
\begin{align}
 K(u,v;a)=g(u)g(v)\bigl\{
 &e^{i\pi(u-v)/2}\Li_{-u-v}(e^{-2\pi ia})\notag\\
 &+e^{-i\pi(u-v)/2}\Li_{-u-v}(e^{2\pi ia})\bigr\}.
 \label{eq:K-polylog}
\end{align}
It is holomorphic near $(u,v)=(0,0)$.

\begin{theorem}[Holomorphic finite-part generating theorem]
\label{thm:generating}
Define the mixed divided difference
\begin{equation}
 R(u,v;a)=\frac{K(u,v;a)-K(u,0;a)-K(0,v;a)+K(0,0;a)}{uv}.
 \label{eq:R-divided}
\end{equation}
The apparent coordinate-axis singularities are removable and
\begin{equation}
 R(u,v;a)=\sum_{m,n\ge0}
       \frac{(-1)^{m+n}}{m!n!}I_{mn}(a)u^mv^n
 \label{eq:R-series}
\end{equation}
in a neighborhood of the origin. In particular,
\begin{align}
 I_{mn}(a)&=(-1)^{m+n}m!n!\coeff{u^{m+1}v^{n+1}}K(u,v;a),
 \label{eq:I-extract}\\
 I_{mn}(a)&=\frac{(-1)^{m+n}}{(m+1)(n+1)}
   \partial_u^{m+1}\partial_v^{n+1}K(0,0;a).
 \label{eq:I-derivative}
\end{align}
\end{theorem}
\begin{proof}
For $\Rea u,\Rea v<0$, Lemma~\ref{lem:convergence} implies
\[
 \int_0^1G(u,x)G(v,x\mathbin\oplus a)\dd x
 =C(1+u,1+v;a)+\frac1{uv}.
\]
Subtracting the endpoint poles therefore gives the meromorphic expression
\begin{equation}
 R=C(1+u,1+v;a)+\frac1{uv}
       +\frac{G(v,a)}u+\frac{G(u,b)}v.
 \label{eq:R-meromorphic}
\end{equation}
Its holomorphy is best proved before any coefficient extraction. Put
\[
 T(u,c)=\frac{1-c^{-u}}u,\qquad T(0,c)=\log c.
\]
A direct local subtraction rewrites \eqref{eq:R-meromorphic} as
\begin{align}
 R={}&\int_0^b\left\{x^{-1-u}[G(v,a+x)-G(v,a)]
                    +G(u,1+x)G(v,a+x)\right\}\dd x\notag\\
 &+\int_0^a\left\{y^{-1-v}[G(u,b+y)-G(u,b)]
                    +G(v,1+y)G(u,b+y)\right\}\dd y\notag\\
 &+G(v,a)T(u,b)+G(u,b)T(v,a).
 \label{eq:R-subtracted}
\end{align}
The differences in square brackets vanish linearly at zero. Consequently
this expression is jointly holomorphic for $u,v$ in a sufficiently small
polydisc about zero; all its spectral derivatives are controlled by
integrable powers times logarithms. Expanding $G$ by
\eqref{eq:stieltjes-definition} gives precisely
\eqref{eq:I-convergent}, coefficient by coefficient, including its upper
endpoint logarithms.

Theorem~\ref{thm:kernel} and $\Gamma(-u)=\Gamma(1-u)/(-u)$ give
$C(1+u,1+v;a)=K(u,v;a)/(uv)$. Moreover the axis values are
\begin{equation}
 K(u,0;a)=-1-uG(u,b),\quad
 K(0,v;a)=-1-vG(v,a),\quad K(0,0;a)=-1.
 \label{eq:K-axes}
\end{equation}
They follow either from \eqref{eq:C-hurwitz} or the Gamma-ratio formula
below. Substitution in \eqref{eq:R-meromorphic} proves
\eqref{eq:R-divided}. Dividing a holomorphic mixed difference by $uv$
removes exactly the axis coefficients, proving \eqref{eq:I-extract}.
Conversion between Taylor coefficients and derivatives gives
\eqref{eq:I-derivative}.
\end{proof}

\begin{remark}[Why a finite part, rather than a product of formal series]
A Fourier expansion at the Stieltjes pole is not an absolutely convergent
series whose factors may be multiplied without qualification.
The proof instead begins in an ordinary convergence region, continues a
specified integral, and matches the continuation to local counterterms.
This matching fixes the finite constant; it is not optional regularization
terminology.
\end{remark}

\section{Gamma-ratio reduction and all-order linear closure}
\label{sec:closure}

Define, with $w=u+v$,
\begin{equation}
 A(u,v)=\frac{\Gamma(1+w)\Gamma(1-v)}{\Gamma(1+u)},
 \qquad B(u,v)=A(v,u).
 \label{eq:A}
\end{equation}
Theorem~\ref{thm:kernel} becomes the particularly useful identity
\begin{equation}
 K(u,v;a)=-uA(u,v)\zeta(1+w,b)-vB(u,v)\zeta(1+w,a).
 \label{eq:K-hurwitz}
\end{equation}
No product of two Hurwitz functions remains on the right.

\begin{lemma}[Removal of the diagonal pole]
\label{lem:Q}
The function
\begin{equation}
 Q(u,v)=\frac{uA(u,v)+vB(u,v)}{u+v}
 \label{eq:Q}
\end{equation}
is holomorphic near zero and
\begin{equation}
 K=-Q-uA(u,v)G(w,b)-vB(u,v)G(w,a).
 \label{eq:K-regular}
\end{equation}
Furthermore
\begin{equation}
 \log A(u,v)=\sum_{k\ge2}\frac{\zeta(k)}k
 \left\{(-1)^k\bigl((u+v)^k-u^k\bigr)+v^k\right\}.
 \label{eq:logA}
\end{equation}
\end{lemma}
\begin{proof}
At $v=-u$, the Gamma ratios $A(u,-u)$ and $B(u,-u)$ both equal one.
The numerator in \eqref{eq:Q} therefore vanishes on $u+v=0$ and is
holomorphically divisible by $u+v$. Splitting off $1/w$ from each Hurwitz
function gives \eqref{eq:K-regular}.
The expansion
$\log\Gamma(1+z)=-\gamma z+\sum_{k\ge2}(-1)^k\zeta(k)z^k/k$
in the defining ratio yields \eqref{eq:logA}. All Euler-constant linear
terms cancel exactly.
\end{proof}

\begin{theorem}[Bilinear closure and the missing-order rule]
\label{thm:closure}
For every $m,n\ge0$, let $N=m+n+1$. There are explicitly computable
coefficients in $\Q[\zeta(2),\ldots,\zeta(N+1)]$ such that
\begin{equation}
 I_{mn}(a)=e_{mn}+\sum_{j=0}^{N}
       \bigl(c_{mn,j}\gamma_j(a)+d_{mn,j}\gamma_j(b)\bigr).
 \label{eq:closure-full}
\end{equation}
The two highest orders satisfy
\begin{equation}
 c_{mn,N}=\frac1{m+1},\qquad d_{mn,N}=\frac1{n+1},\qquad
 c_{mn,N-1}=d_{mn,N-1}=0.
 \label{eq:top-coefficients}
\end{equation}
The symmetry is $c_{mn,j}=d_{nm,j}$ and $e_{mn}=e_{nm}$.
Assigning formal weight $k$ to $\zeta(k)$, the coefficients
$c_{mn,j},d_{mn,j}$ are homogeneous of weight $N-j$, and $e_{mn}$ is
homogeneous of weight $N+1$. Zero coefficients obey the same convention.
\end{theorem}
\begin{proof}
Set $D=N+1=m+n+2$ and $f=(-1)^{m+n}m!n!$.
Let $A_d$ be the homogeneous part of total degree $d$ in $A$, and put
$B_d(u,v)=A_d(v,u)$. Equation~\eqref{eq:K-regular} gives the finite
coefficient formulas
\begin{align}
 e_{mn}&=-f\coeff{u^{m+1}v^{n+1}}
              \frac{uA_D+vB_D}{u+v},\label{eq:coefficient-e}\\
 c_{mn,j}&=-\frac{f(-1)^j}{j!}
       \coeff{u^{m+1}v^{n+1}}vB_{N-j}(u,v)(u+v)^j,\label{eq:coefficient-c}\\
 d_{mn,j}&=-\frac{f(-1)^j}{j!}
       \coeff{u^{m+1}v^{n+1}}uA_{N-j}(u,v)(u+v)^j.
       \label{eq:coefficient-d}
\end{align}
Each numerator in \eqref{eq:coefficient-e} is polynomially divisible by
$u+v$, by homogeneous decomposition of Lemma~\ref{lem:Q}.
Every term in \eqref{eq:logA} has spectral degree and zeta weight equal
to its index. This proves the finite coefficient algebra and weight
statements.

For $j=N$, use $A_0=B_0=1$. The binomial coefficient in
$v(u+v)^N$ gives
\[
 -\frac{(-1)^{m+n+N}m!n!}{N!}\binom{N}{m+1}
 =\frac1{m+1},
\]
and the other coefficient is $1/(n+1)$. For $j=N-1$, both expressions
contain $A_1$ or $B_1$, which are zero because \eqref{eq:logA} has no
linear term. The symmetry follows by swapping $u,v$ and the two
Stieltjes arguments.
\end{proof}

The theorem does not treat positive zeta values as algebraically
independent. Formal homogeneity describes the explicit expressions
produced by the Gamma expansion. Classical relations among even zeta
values may be applied afterward, or the unsimplified expressions may be
retained to preserve an auditable coefficient convention.

\subsection{A terminating exact coefficient algorithm}
Write $\ell_k$ for the degree-$k$ summand of \eqref{eq:logA}.
The homogeneous exponential satisfies
\begin{equation}
 A_0=1,\qquad A_1=0,\qquad
 A_d=\frac1d\sum_{k=2}^d k\ell_k A_{d-k}\quad(d\ge2).
 \label{eq:recurrence}
\end{equation}
To prove it, insert an auxiliary variable $r$ into
$\exp(\sum_{k\ge2}r^k\ell_k)$, differentiate in $r$, and compare
coefficients. For a fixed pair $(m,n)$, only degrees $d\le m+n+2$ are
needed. Thus the calculation terminates by rational polynomial
arithmetic, without integer-relation searches or numerical constants.
The program \file{code/exact_coefficients.py} implements exactly
\eqref{eq:recurrence} and \eqref{eq:coefficient-e}--\eqref{eq:coefficient-d}.

\section{Explicit identities and a convergent digamma formula}
\label{sec:examples}

For compactness in this section only, write
$a_j=\gamma_j(a)$, $b_j=\gamma_j(b)$ and $z_k=\zeta(k)$.
The first identities are
\begin{align}
 I_{00}&=a_1+b_1-2z_2,\label{eq:I00}\\
 I_{10}&=\tfrac12a_2+b_2+z_2(a_0+b_0)-z_3,\label{eq:I10}\\
 I_{11}&=\tfrac12(a_3+b_3)+2z_2(a_1+b_1)
                       +z_3(a_0+b_0)-z_2^2-z_4,\label{eq:I11}\\
 I_{20}&=\tfrac13a_3+b_3+2z_2(a_1+b_1)+2z_3b_0-z_2^2-3z_4,
                       \label{eq:I20}\\
 I_{21}&=\tfrac13a_4+\tfrac12b_4+3z_2(a_2+b_2)
                  +2z_3a_1+4z_3b_1\notag\\
       &\quad +(2z_2^2+2z_4)a_0+(z_2^2+3z_4)b_0
                    -4z_2z_3-4z_5,\label{eq:I21}\\
 I_{30}&=\tfrac14a_4+b_4+3z_2(a_2+b_2)+6z_3b_1
                    +(3z_2^2+3z_4)a_0+6z_4b_0-6z_5.
                    \label{eq:I30}
\end{align}
All formulas with the indices reversed follow by $a\leftrightarrow b$.
For example, the lack of $a_2,b_2$ in $I_{11}$ is the universal
missing-order rule, not an accidental low-weight simplification.

\begin{corollary}[Ordinary convergent digamma integrals]
\label{cor:digamma}
For $0<a<1$, $b=1-a$,
\begin{align}
 &\int_0^b\left\{\psi(x)\psi(x+a)+\frac{\psi(a)}x\right\}\dd x
 +\int_0^a\left\{\psi(y)\psi(y+b)+\frac{\psi(b)}y\right\}\dd y\notag\\
 &\hspace{10mm}=\gamma_1(a)+\gamma_1(b)-\frac{\pi^2}{3}
                       +\psi(a)\log b+\psi(b)\log a.
 \label{eq:digamma-convergent}
\end{align}
Both displayed integrals converge absolutely.
\end{corollary}
\begin{proof}
Insert $\gamma_0=-\psi$ into \eqref{eq:I-convergent} and use
\eqref{eq:I00}. Near zero, $\psi(x)=-1/x-\gamma+O(x)$, so each
subtracted integrand is bounded. This also verifies the convergence
directly, without relying on finite-part notation.
\end{proof}

At $a=1/2$, the distribution identity gives
\[
 I_{00}(1/2)=2\gamma_1-4\gamma\log2-2\log^22-\frac{\pi^2}{3}.
\]
This follows from $\gamma_1(1/2)=\gamma_1-2\gamma\log2-\log^22$.
The defining quadratic integral has been replaced by a short linear
expression in the ordinary Stieltjes constants and elementary constants.

\section{Coincident endpoints and canonical collision subtraction}
\label{sec:coincident}

A translated correlation has two separated logarithmic endpoint
divergences. At zero translation they merge into an additional
$x^{-2}$ singularity. We now show that the same coefficient table still
applies, but only after the merged singularity is subtracted explicitly.

For an integer $r\ge0$, put
\begin{equation}
 E_r(L)=\sum_{k=0}^r\frac{r!}{k!}L^k.
 \label{eq:E-polynomial}
\end{equation}
The identity $E_r-E_r'=L^r$ implies
\begin{equation}
 \int_\epsilon^1\frac{\log^r x}{x^2}\dd x
 =\frac{E_r(\log\epsilon)}{\epsilon}-r!.
 \label{eq:double-endpoint}
\end{equation}

\begin{definition}[Coincident-endpoint finite part]
\label{def:J}
Write $\gamma_j=\gamma_j(1)$ and define
\begin{align}
 J_{mn}=\lim_{\epsilon\downarrow0}\biggl\{
 &\int_\epsilon^1\gamma_m(x)\gamma_n(x)\dd x
       -\frac{E_{m+n}(\log\epsilon)}\epsilon\notag\\
 &+\frac{\gamma_n}{m+1}\log^{m+1}\epsilon
       +\frac{\gamma_m}{n+1}\log^{n+1}\epsilon\biggr\}.
 \label{eq:J-cutoff}
\end{align}
\end{definition}

\begin{theorem}[Coincident closure]
\label{thm:coincident}
The finite part \eqref{eq:J-cutoff} exists. With exactly the coefficients
of Theorem~\ref{thm:closure},
\begin{equation}
 J_{mn}=e_{mn}+\sum_{j=0}^{m+n+1}(c_{mn,j}+d_{mn,j})\gamma_j.
 \label{eq:J-closure}
\end{equation}
Equivalently it is the same mixed coefficient of
\begin{equation}
 K_0(u,v)=-(uA(u,v)+vB(u,v))\zeta(1+u+v).
 \label{eq:K0}
\end{equation}
In particular,
\begin{equation}
 J_{00}=2\gamma_1-\frac{\pi^2}{3}.
 \label{eq:J00}
\end{equation}
\end{theorem}
\begin{proof}
By \eqref{eq:stieltjes-shift}, the cutoff in Definition~\ref{def:J}
converges to the absolutely convergent expression
\begin{align}
 J_{mn}={}&-(m+n)!\notag\\
 &+\int_0^1\biggl\{
 \frac{\log^m x}{x}[\gamma_n(1+x)-\gamma_n]
 +\frac{\log^n x}{x}[\gamma_m(1+x)-\gamma_m]\notag\\
 &\hspace{35mm}+\gamma_m(1+x)\gamma_n(1+x)\biggr\}\dd x.
 \label{eq:J-convergent}
\end{align}
The two bracketed differences vanish linearly at zero, and
\eqref{eq:double-endpoint} supplies the constant $-(m+n)!$.

To identify the coefficients, start where the ordinary coincident
Hurwitz integral converges, for example $\Rea s,\Rea t<0$.
Fourier orthogonality gives
\[
 C_0(s,t)=2\Gamma(1-s)\Gamma(1-t)(2\pi)^{s+t-2}
       \cos\!\left(\frac{\pi(s-t)}2\right)\zeta(2-s-t).
\]
The Riemann functional equation, equivalently the specialization of the
Fourier calculation in Theorem~\ref{thm:kernel}, gives
$uvC_0(1+u,1+v)=K_0(u,v)$. Meromorphic continuation is justified by
subtracting the local terms as above.

More explicitly, the holomorphic generating expression near zero is
\begin{align*}
 R_0(u,v)={}&-\frac1{1+u+v}\\
 &+\int_0^1\bigl\{
 x^{-1-u}[G(v,1+x)-G(v,1)]
 +x^{-1-v}[G(u,1+x)-G(u,1)]\\
 &\hspace{39mm}+G(u,1+x)G(v,1+x)\bigr\}\dd x.
\end{align*}
The first term is the analytic continuation of
$\int_0^1x^{-2-u-v}\dd x$, not a discarded divergence. The remaining
integrals converge locally uniformly near zero. On a starting convergence
region, and hence by continuation,
\[
 R_0=\frac{K_0+1+uG(u,1)+vG(v,1)}{uv}.
\]
Its Taylor coefficients are \eqref{eq:J-convergent}. Decomposing
\eqref{eq:K0} as $-Q-(uA+vB)G(u+v,1)$ now gives
\eqref{eq:J-closure} by the identical finite algebra used earlier.
\end{proof}

For the first case, an ordinary integral version is worth recording:
\begin{equation}
 \int_0^1\left\{\psi(1+x)^2
       -\frac{2[\psi(1+x)+\gamma]}x\right\}\dd x
 =1+2\gamma_1-\frac{\pi^2}{3}.
 \label{eq:digamma-square-convergent}
\end{equation}
The apparently singular numerator in the second term vanishes at zero.
There is no positivity contradiction in the negative value of $J_{00}$:
it is a finite part of a positive divergent square, not its ordinary
nonnegative integral.

\subsection{An exact convergent expansion at collision}
The relationship between the separated and coincident finite parts can
be expressed by an equality of convergent series, not merely an
asymptotic prescription.

\begin{theorem}[Collision polynomial and convergent regular part]
\label{thm:collision}
For $m,n\ge0$ define a polynomial
\begin{equation}
 \mathcal P_{mn}(L)=(-1)^{m+n+1}m!n!
       \coeff{u^{m+1}v^n}B(u,v)e^{-(u+v)L}.
 \label{eq:collision-polynomial}
\end{equation}
It has degree $m+n+1$, leading coefficient $1/(m+1)$, and coefficients
in positive integer zeta values. There is an exact expansion
\begin{equation}
 I_{mn}(a)=\frac{\mathcal P_{mn}(\log a)}a
                +J_{mn}+\sum_{k\ge1}h_{mn,k}a^k,
 \qquad 0<a<1,
 \label{eq:collision-series}
\end{equation}
whose regular power series converges locally uniformly for $|a|<1$.
Its coefficients are
\begin{align}
 h_{mn,k}={}&\frac{(-1)^{m+n+1}m!n!}{k!}
 \coeff{u^{m+1}v^{n+1}}\notag\\
 &\quad\bigl(uA(u,v)+(-1)^kvB(u,v)\bigr)
               (1+u+v)_k\zeta(1+u+v+k).
 \label{eq:collision-coefficients}
\end{align}
For $m=n$, all odd-indexed $h_{mm,k}$ vanish. In particular,
\begin{equation}
 \lim_{a\downarrow0}\left\{I_{mn}(a)
             -\frac{\mathcal P_{mn}(\log a)}a\right\}=J_{mn}.
 \label{eq:collision-limit}
\end{equation}
\end{theorem}
\begin{proof}
The classical parameter Taylor series at one gives, locally uniformly
in the spectral variable,
\[
 \zeta(1+w,1-a)=\sum_{k\ge0}\frac{(1+w)_k}{k!}
                      \zeta(1+w+k)a^k,
\]
while the shift law gives
\[
 \zeta(1+w,a)=a^{-1-w}
        +\sum_{k\ge0}\frac{(1+w)_k}{k!}\zeta(1+w+k)(-a)^k.
\]
These are meromorphic in $w$ at zero, with the $k=0$ pole explicitly
visible. Substitution in \eqref{eq:K-hurwitz} yields
\begin{align}
 K(u,v;a)={}&-vB(u,v)a^{-1-w}\notag\\
 &-\sum_{k\ge0}\frac{uA(u,v)+(-1)^kvB(u,v)}{k!}
                      (1+w)_k\zeta(1+w+k)a^k.
 \label{eq:K-collision}
\end{align}
The $k=0$ summand is exactly $K_0$; its apparent pole at $w=0$ is
removable. The remaining series is holomorphic for small $u,v$ and
$|a|<1$, uniformly on compact subsets. For example, on a fixed small
spectral polydisc its terms and spectral derivatives grow at most
polynomially in $k$ times powers of $\log k$, whereas $|a|^k$ decays
geometrically on compact subdiscs. Coefficient extraction is therefore
legitimate. The singular first term gives
\eqref{eq:collision-polynomial}, the constant gives
Theorem~\ref{thm:coincident}, and the other terms give
\eqref{eq:collision-coefficients}.

The highest power of $L$ comes from $B_0=1$ and has degree $m+n+1$;
a binomial coefficient gives its stated leading coefficient.
For $m=n$ and odd $k$, the factor
$uA-vB$ is antisymmetric under $u\leftrightarrow v$, while every other
factor in \eqref{eq:collision-coefficients} is symmetric. Its diagonal
mixed coefficient is zero. Finally the regular series has zero limit
at $a=0$, proving \eqref{eq:collision-limit}.
\end{proof}

The first collision polynomials are
\begin{equation}
 \mathcal P_{00}(L)=L,\qquad
 \mathcal P_{10}(L)=\frac{L^2}{2}+\zeta(2),\qquad
 \mathcal P_{01}(L)=L^2+\zeta(2).
 \label{eq:collision-examples}
\end{equation}
The asymmetry of the last two is required: the factor that becomes
singular at $a=0$ has changed its index. In particular,
\[
 I_{00}(a)-\frac{\log a}{a}\longrightarrow
                   2\gamma_1-\frac{\pi^2}{3}.
\]
There is no arbitrary additional finite constant in this limit. It is
fixed by the initial endpoint subtraction and agrees with the coincident
prescription.

\section{An ordinary log-Gamma autocorrelation identity}
\label{sec:loggamma}

The same two-parameter kernel has a regular spectral specialization at
zero. Unlike the Stieltjes correlations, the following integral needs no
finite-part interpretation:
\begin{align}
 F(a)={}&\int_0^b\log\Gamma(x)\log\Gamma(x+a)\dd x\notag\\
       &+\int_0^a\log\Gamma(b+y)\log\Gamma(y)\dd y,
 \qquad 0<a<1.
 \label{eq:F-definition}
\end{align}
The endpoint singularities are logarithmic and integrable. Define
$L=\log(2\pi)$, $\alpha=\gamma+L$ and
$B_2(a)=a^2-a+1/6$.

\begin{theorem}[Zeta and polylogarithm forms of the autocorrelation]
\label{thm:loggamma}
For $0<a<1$,
\begin{align}
 F(a)={}&\frac{L^2}{4}+(1+\zeta(2))B_2(a)
        -\zeta'(-1,a)-\zeta'(-1,b)\notag\\
        &\hspace{29mm}-\frac12\bigl(\zeta''(-1,a)+\zeta''(-1,b)\bigr).
 \label{eq:F-hurwitz}
\end{align}
With $z=e^{2\pi ia}$, an equivalent identity is
\begin{equation}
 F(a)=\frac{L^2}{4}+\frac{1}{2\pi^2}
 \left\{\left(\alpha^2+\frac{\pi^2}{4}\right)\Rea\Li_2(z)
       -2\alpha\Rea\LiJ{1}{z}+\Rea\LiJ{2}{z}\right\}.
 \label{eq:F-polylog}
\end{equation}
Both formulas extend continuously to $a=0$ and $a=1$.
\end{theorem}
\begin{proof}
Lerch's formula gives
$\zeta'(0,x)=\log\Gamma(x)-L/2$ and its integral over $[0,1]$ is
zero. Differentiating \eqref{eq:C-definition} once in each spectral
variable at zero is legitimate by Lemma~\ref{lem:convergence}; hence
\begin{equation}
 F(a)=L^2/4+\partial_s\partial_tC(0,0;a).
 \label{eq:F-derivative}
\end{equation}
For the first expression let $w=s+t$ and use the same Gamma ratio $A$
as in \eqref{eq:A}, now in variables $s,t$. Equation~\eqref{eq:C-hurwitz}
rewrites as
\begin{equation}
 C(s,t;a)=\frac{sA(s,t)\zeta(w-1,b)+tB(s,t)\zeta(w-1,a)}{w(w-1)}.
 \label{eq:C-near-zero}
\end{equation}
Since $\zeta(-1,a)=\zeta(-1,b)=-B_2(a)/2$, the $w=0$
singularity is removable. Also
$A(s,t)=1+\zeta(2)t(s+t)+O((|s|+|t|)^3)$.
Expanding the numerator through total degree three and extracting $st$
after division gives
\begin{align*}
 \partial_s\partial_t C(0,0;a)
 ={}&-2(1+\zeta(2))\zeta(-1,a)\\
    &-\zeta'(-1,a)-\zeta'(-1,b)
      -\tfrac12\bigl(\zeta''(-1,a)+\zeta''(-1,b)\bigr).
\end{align*}
This proves \eqref{eq:F-hurwitz}.

Alternatively, differentiate the Fourier-coefficient identity
\eqref{eq:fourier-coefficients} at $s=0$. This is justified by the
endpoint domination in Lemma~\ref{lem:convergence}, applied to each
coefficient integral. For $n>0$ the coefficient of the centered
log-Gamma function is
\[
 -\frac{i}{2\pi n}\left(\alpha+\log n+\frac{i\pi}{2}\right).
\]
Its product with the coefficient at $-n$ is
$((\alpha+\log n)^2+\pi^2/4)/(4\pi^2n^2)$.
The centered log-Gamma function is square-integrable, so the bilinear
Parseval identity for its translation consequently gives
\[
 F(a)-L^2/4=\frac1{2\pi^2}\sum_{n\ge1}
 \frac{(\alpha+\log n)^2+\pi^2/4}{n^2}\cos(2\pi na).
\]
The series is absolutely and uniformly convergent. Expanding the square
and using the spectral derivative series of $\Li_s$ at $s=2$ gives
\eqref{eq:F-polylog}. Uniform convergence establishes continuity on the
closed interval. Equivalently, log-Gamma is square-integrable and its
circular autocorrelation is continuous.
\end{proof}

\subsection{Half- and quarter-period evaluations}
Write $\ell=\log2$ and $Z_j=\zeta^{(j)}(-1)$. Differentiating the elementary
Hurwitz distribution identities at $-1$ gives
\begin{align}
 F(1/2)={}&\frac{L^2}{4}+(1-\ell)Z_1+\frac{Z_2}{2}
       +\frac{\ell-1-\zeta(2)}{12}+\frac{\ell^2}{24},
 \label{eq:F-half}\\
 F(1/4)={}&\frac{L^2}{4}+\frac{Z_1}{4}+\frac{Z_2}{8}
       +\frac{\ell^2-1-\zeta(2)}{48}.
 \label{eq:F-quarter}
\end{align}
For the first identity use $\zeta(s,1/2)=(2^s-1)\zeta(s)$.
For the second use
\[
 \zeta(s,1/4)+\zeta(s,3/4)=(4^s-2^s)\zeta(s).
\]
Thus these translated products require only two Riemann-zeta jets at
$-1$, not an unresolved numerical relation between a large set of
Gamma values.

The zero-shift limit reads
\begin{equation}
 F(0)=\frac{L^2}{4}+\frac{1+\zeta(2)}6
                          -2\zeta'(-1)-\zeta''(-1).
 \label{eq:F-zero}
\end{equation}
Differentiating the Riemann functional equation converts it to the
classical quadratic log-Gamma moment recorded in the manuscript
\cite{repo-integration}:
\[
 \int_0^1\log^2\Gamma(x)\dd x
 =\frac{\gamma^2}{12}+\frac{\pi^2}{48}+\frac{\gamma L}{6}
  +\frac{L^2}{3}-\frac{(\gamma+L)\zeta'(2)}{\pi^2}
  +\frac{\zeta''(2)}{2\pi^2}.
\]
This is a compatibility check, not a claim to rediscover the classical
moment as a new theorem.

\section{Derivatives, polygamma and exact antiderivatives}
\label{sec:primitives}

The closure theorem is useful not only because it eliminates a product.
It also converts differentiation and integration of the entire family
into standard operations on one Hurwitz function.

\begin{proposition}[Parameter-derivative dictionary]
\label{prop:derivative}
For integers $r\ge1$, $n\ge0$, and $a>0$,
\begin{equation}
 \gamma_n^{(r)}(a)=(-1)^{r+n}n!\coeff{u^n}
                  (1+u)_r\zeta(r+1+u,a).
 \label{eq:gamma-derivatives}
\end{equation}
In particular,
\begin{equation}
 \gamma_1^{(r)}(a)=(-1)^{r+1}r!
                   \bigl(\zeta'(r+1,a)+H_r\zeta(r+1,a)\bigr).
 \label{eq:gamma1-derivatives}
\end{equation}
\end{proposition}
\begin{proof}
Iterate the Hurwitz parameter derivative identity to obtain
$\partial_a^r\zeta(1+u,a)=(-1)^r(1+u)_r\zeta(r+1+u,a)$.
The derivative of $1/u$ in $a$ is zero. Coefficient comparison in
\eqref{eq:stieltjes-definition} gives \eqref{eq:gamma-derivatives}.
The first two terms of $(1+u)_r=r!(1+H_ru+O(u^2))$ give
\eqref{eq:gamma1-derivatives}.
\end{proof}

For instance,
\begin{align}
 \frac{\mathrm d^r}{\mathrm da^r}I_{00}(a)
 =(-1)^{r+1}r!\bigl\{&\zeta'(r+1,a)+H_r\zeta(r+1,a)\notag\\
 &+(-1)^r[\zeta'(r+1,b)+H_r\zeta(r+1,b)]\bigr\}.
 \label{eq:I00-derivatives}
\end{align}
Ordinary polygamma values enter through
$\psi^{(r)}(a)=(-1)^{r+1}r!\zeta(r+1,a)$.
Equations \eqref{eq:gamma-derivatives} and \eqref{eq:closure-full} give
all higher derivatives of every $I_{mn}$ without differentiating an
ill-defined product at a singular endpoint.

\begin{proposition}[A closed primitive of every correlation]
\label{prop:primitive}
Suppose $I_{mn}(a)$ is written as in \eqref{eq:closure-full}.
A primitive on $0<a<1$ is
\begin{equation}
 e_{mn}a+\sum_{j=0}^{m+n+1}\frac{(-1)^{j+1}}{j+1}
 \left\{c_{mn,j}\zeta^{(j+1)}(0,a)
              -d_{mn,j}\zeta^{(j+1)}(0,1-a)\right\}.
 \label{eq:general-primitive}
\end{equation}
In particular,
\begin{equation}
 \int I_{00}(a)\dd a
 =\frac{\zeta''(0,a)-\zeta''(0,1-a)}2
                  -\frac{\pi^2}{3}a+\text{constant}.
 \label{eq:I00-primitive}
\end{equation}
\end{proposition}
\begin{proof}
Expand $\partial_a\zeta(s,a)=-s\zeta(s+1,a)$ at $s=0$.
Comparing coefficients gives
\[
 \partial_a\zeta^{(j+1)}(0,a)=(-1)^{j+1}(j+1)\gamma_j(a).
\]
This is the classical antiderivative ladder already used in the
repository \cite{repo-integration}. Apply it termwise to
\eqref{eq:general-primitive}; the reflected argument contributes a minus
sign by the chain rule. Equation~\eqref{eq:I00-primitive} is the first
specialization.
\end{proof}

Repeated primitives can also be written with no recursively undefined
functions. If $r\ge1$ and $j\ge0$, then
\begin{equation}
 \left.\partial_s^j
 \left\{\frac{\zeta(s-r,a)}{\prod_{k=1}^r(k-s)}\right\}\right|_{s=0}
 \label{eq:repeated-primitive}
\end{equation}
is an $r$-fold primitive of $\zeta^{(j)}(0,a)$.
Indeed $r$ parameter derivatives of the numerator multiply it by
$\prod_{k=1}^r(k-s)$ and restore $\zeta(s,a)$.
The integration polynomial of degree at most $r-1$ can be fixed by any
specified basepoint normalization. At $j=1$, these formulas are one way
to represent the negative-polygamma ladder; they do not silently change
the zero-based or balanced conventions in the existing manuscript.

\section{Rational grids and Dirichlet-character coordinates}
\label{sec:arithmetic}

The finite linear closure interacts particularly well with classical
Hurwitz distribution. For integers $q\ge2$ put
\begin{equation}
 T_j(q)=\sum_{r=1}^{q-1}\gamma_j(r/q).
 \label{eq:T-definition}
\end{equation}

\begin{proposition}[Exact rational-grid traces]
\label{prop:trace}
For $j\ge0$,
\begin{equation}
 T_j(q)=(q-1)\gamma_j+
 q\sum_{k=1}^j\binom{j}{k}(-\log q)^k\gamma_{j-k}
       +\frac{(-1)^j q}{j+1}(\log q)^{j+1}.
 \label{eq:T-formula}
\end{equation}
Consequently
\begin{equation}
 \sum_{r=1}^{q-1}I_{mn}(r/q)
 =(q-1)e_{mn}+\sum_{j=0}^{m+n+1}
                    (c_{mn,j}+d_{mn,j})T_j(q).
 \label{eq:I-trace}
\end{equation}
The first case is
\begin{equation}
 \sum_{r=1}^{q-1}I_{00}(r/q)
 =2(q-1)(\gamma_1-\zeta(2))-2q\gamma\log q-q\log^2q.
 \label{eq:I00-trace}
\end{equation}
\end{proposition}
\begin{proof}
Sum the Hurwitz distribution identity and remove the $r=q$ term:
\[
 \sum_{r=1}^{q-1}\zeta(1+u,r/q)
                         =(q^{1+u}-1)\zeta(1+u).
\]
The residue on both sides is $(q-1)/u$. Expanding the regular part and
multiplying its $u^j$ coefficient by $(-1)^jj!$ gives
\eqref{eq:T-formula}; the last term comes from the pole times the
exponential $q^u$, and must be retained. Summing the closure formula and
permuting $r\leftrightarrow q-r$ gives \eqref{eq:I-trace}.
\end{proof}

Let $\chi$ be any nonprincipal Dirichlet character modulo $q$, not
necessarily primitive, and define
$W_j(\chi)=\sum_{r=1}^{q-1}\chi(r)\gamma_j(r/q)$.
Our convention uses $\chi$, not its complex conjugate, in this sum.

\begin{proposition}[Character resolution]
\label{prop:character}
For $j\ge0$,
\begin{equation}
 W_j(\chi)=(-1)^jq\sum_{k=0}^j\binom{j}{k}
                (\log q)^{j-k}L^{(k)}(1,\chi).
 \label{eq:W-formula}
\end{equation}
Moreover
\begin{equation}
 \sum_{r=1}^{q-1}\chi(r)I_{mn}(r/q)
 =\sum_{j=0}^{m+n+1}
           \bigl(c_{mn,j}+\chi(-1)d_{mn,j}\bigr)W_j(\chi).
 \label{eq:character-correlation}
\end{equation}
In particular the character trace of $I_{mm}$ vanishes for every odd
character.
\end{proposition}
\begin{proof}
The identity
$\sum_r\chi(r)\zeta(s,r/q)=q^sL(s,\chi)$ has no pole at $s=1$,
because $\sum_r\chi(r)=0$. Differentiating $j$ times there gives
\eqref{eq:W-formula}. The constant in the closure formula sums to zero,
and the reflected Stieltjes term transforms by
$\chi(q-r)=\chi(-1)\chi(r)$. This proves
\eqref{eq:character-correlation}. For $m=n$, the two coefficients are
equal and their odd-character contribution cancels.
\end{proof}

These statements reduce arithmetic traces of nonlinear-looking
integrals to a finite spectral jet of a single Dirichlet $L$-function.
They do not prove that the surviving even-character derivative
coordinates are independent of other periods. No such independence
assumption is used in their derivation.

\section{Computational reproduction and evidence separation}
\label{sec:verification}

The proof of Theorem~\ref{thm:closure} is uniform in the indices.
Computations serve a different purpose: detecting coefficient, phase,
normalization and implementation mistakes. The supplied programs retain
this separation in their module documentation and result files.

\subsection{Exact algebra}
The coefficient engine generates all $45$ ordered pairs $(m,n)$ with
$m+n\le8$. Zeta values are formal symbols during this calculation.
The run recorded in \file{results/exact_verification.json} checks
$272$ exact assertions: polynomial divisibility is enforced during
coefficient construction; the counted assertions verify reflection
symmetry, the two highest nonzero Stieltjes coefficients, the missing
next order, the leading collision coefficient, and two explicitly
specified base identities. The result is \texttt{PASS}.

The full data in \file{results/coefficients.json} store the constant,
the coefficients of $\gamma_j(a)$ and $\gamma_j(1-a)$, and the collision
polynomial for every pair. The companion \file{results/identities.txt}
is a human-readable rendering. No transcendental equality is inferred
from rational rank computations, and the symbolic checks are not a
formalization of the convergence proofs.

\subsection{Independent numerical diagnostics}
With Python 3.13.5, mpmath 1.3.0 and $45$ decimal working digits,
\file{code/verify_numeric.py} passed $19$ checks. They compare ordinary
or explicitly subtracted quadrature with the closed formulas, compare
polylogarithm and Hurwitz kernels at noninteger real and complex orders,
check rational traces, and differentiate a closed antiderivative.
The largest recorded absolute residual in this run is approximately
$1.06\times10^{-44}$.

A separate program, \file{code/verify_stieltjes_pairs.py}, evaluates the
regularized products $I_{10},I_{11},I_{20}$ at $a=1/2$ by quadrature.
It evaluates the necessary Stieltjes functions through a $140$-term
parameter Taylor expansion at one, using Hurwitz spectral derivatives
for the coefficients, rather than evaluating the closed correlation
formula inside its own integrand. The recorded comparisons are
\begin{center}
\begin{tabular}{@{}lrl@{}}
\toprule
Quantity & Quadrature value (rounded) & Absolute residual\\
\midrule
$I_{10}(1/2)$ & $6.71092887447221586072$ & $3.98\times10^{-41}$\\
$I_{11}(1/2)$ & $-8.64046177754241055451$ & $1.61\times10^{-40}$\\
$I_{20}(1/2)$ & $-11.02758300286848051739$ & $2.34\times10^{-40}$\\
\bottomrule
\end{tabular}
\end{center}
The Taylor truncation and quadrature errors are not enclosed by directed
interval arithmetic. The table is therefore numerical evidence, not
certified error bounds. The combined numerical count is $22$, separate
from the $272$ exact symbolic checks.

A negative control omits the term $-2\zeta(2)$ from $I_{00}$ and verifies
that the resulting discrepancy is $2\zeta(2)$. This tests the particular
finite constant that is easy to lose when endpoint subtraction is
replaced by an informal cancellation of divergent Fourier terms.

\subsection{A reproducible order-differentiation canary}
\label{sec:canary}
At integer spectral orders, differentiating a polylogarithm numerical
routine with default finite-difference settings can lose substantially
more digits than direct function evaluation suggests. The exact identity
\begin{equation}
 \Rea\Li_s(i)=2^{-s}(2^{1-s}-1)\zeta(s)
 \label{eq:cyclotomic-canary}
\end{equation}
for real $s$ near $2$ supplies a simple test. It follows by grouping
the absolutely convergent series by even indices, and then continues
in $s$.

In the installed mpmath 1.3.0, at $40$ working digits and with the exact
argument \texttt{mp.mpc(0,1)}, applying the default \texttt{mp.diff}
to the left and right sides of \eqref{eq:cyclotomic-canary} gives zero
recorded discrepancy for the values, but absolute discrepancies
\[
 2.33945455584\times10^{-12}\quad\hbox{and}\quad
 2.44970539683\times10^{-16}
\]
for the first and second spectral derivatives, respectively.
The minimal reproducer and full observed values are supplied in
\file{code/polylog_derivative_canary.py} and its JSON output.

This is a version- and method-specific observation, not a counterexample
to \eqref{eq:cyclotomic-canary}, nor a claim that every polylogarithm
implementation has the same behavior. The validating log-Gamma tests
therefore use exact cyclotomic reduction before spectral differentiation
at $i$ and $-1$. The raw polylogarithm kernel is also checked separately
at noninteger spectral parameters. Future numerical implementations may
remove the canary discrepancy; the reproducer records it rather than
requiring it to persist.

\section{Audit, normalization safeguards and integration scope}
\label{sec:audit}

\paragraph{Finite part versus ordinary integral.}
The new formulas must be integrated with Definition~\ref{def:finitepart},
not with a bare integral sign. Counterterms include both endpoint
coefficients and the correct powers of logarithms. At collision, the
extra $\epsilon^{-1}E_{m+n}(\log\epsilon)$ term is indispensable.
For users who prefer ordinary integrals,
\eqref{eq:I-convergent}, \eqref{eq:digamma-convergent} and
\eqref{eq:J-convergent} eliminate finite-part notation entirely.

\paragraph{Analytic continuation and diagonal cancellation.}
Gamma poles in \eqref{eq:C-hurwitz} and the apparent pole at $u+v=0$
in \eqref{eq:K-hurwitz} cancel only in the combined expression.
Expanding or evaluating the individual terms at a pole is invalid.
The regular expression \eqref{eq:K-regular} or exact homogeneous
coefficients should be used near the intersection of the axes.
The limit $(1-c^{-u})/u\to\log c$ also fixes an easily reversed sign in
subtraction code.

\paragraph{Differentiating singular products.}
The parameter derivatives in Section~\ref{sec:primitives} differentiate
a proved closed expression. We do not assert that differentiation
commutes with an unmodified finite-part product integral. Moving the
break point, differentiating endpoint values, and differentiating
counterterms can all contribute. A separate commutation theorem would
be required for that claim.

\paragraph{Repository status and existing corrections.}
The inspected integration material already distinguishes classical
identities from experimental searches, and its denominator-five
log-Gamma correction retains the required Dirichlet derivative
\cite{repo-integration}. We do not replace that term by an unsupported
Clausen-only reduction. The current exact $S_4$ proof must retain its
proved status \cite{repo-s4}. Neither the $S_6$ problem nor period
independence has been settled by the present bilinear calculus.

\paragraph{Computational finding, not an invented source erratum.}
The observed polylogarithm order-derivative accuracy loss is a real
reproducible finding in the delivery environment. It motivates a
regression canary for future repository computations. We have not
established that a particular existing repository result relied on the
failing default differentiation path, so no existing theorem is
retracted on that basis.

\paragraph{Review boundaries.}
The read of \file{main} at the end of the review returned commit
\file{e1294e75708b84986bfdaf103892ac9ac78a0580}.
Individual selected file blob identifiers are recorded separately in
\file{notes/PROVENANCE.json}; the earlier reads were not all pinned to
that commit. Incoming ZIPs were inventoried but not unpacked or audited.
The package proposes a new report and a namespaced integration section,
without modifying any live repository file. This respects the incoming
archive distinction documented by the repository \cite{repo-incoming}.

\section{Further research questions}
\label{sec:questions}

\subsection{A trilinear analogue and its genuine new coordinates}
For three distinct translations, define a product finite part by
subtracting the local singular terms at all three endpoints. Does its
spectral generating function reduce to finitely many Tornheim-type or
colored double-zeta jets, and what is a useful exact coefficient normal
form? Unlike the bilinear Fourier calculation, a trilinear integral
retains a nontrivial additive relation among three frequencies. A
linear closure theorem analogous to Theorem~\ref{thm:closure} should not
be presumed without proving those remaining relations.

\subsection{Unequal integer dilations}
Replace the translated factor by a periodic function of $qx+a$, or
allow different positive integer dilations in the two factors. Fourier
orthogonality imposes a divisibility condition instead of equality of
frequencies. Determine a finite distribution formula for the resulting
finite parts, including the endpoints introduced by the dilation. This
could connect the Gamma-ratio coefficient algebra directly to the
manuscript's rational-grid modules.

\subsection{Parameter differentiation with moving counterterms}
Prove an explicit commutation law between differentiation in $a$ and the
finite-part operation. The desired theorem should list every endpoint
term and yield independent integral identities involving higher
polygamma and Stieltjes derivatives, rather than merely restating the
closed parameter derivative in Section~\ref{sec:primitives}.

\subsection{Higher integral spectral resonances}
The expansions here are centered at spectral pairs $(1,1)$ and $(0,0)$.
At other integer pairs, Gamma factors and Hurwitz poles can meet in
more complicated ways. Classify the corresponding removable combinations
and determine which products reduce to finite single-function jets.
The first useful examples involve a Stieltjes factor paired with a
polygamma or an explicitly normalized negative-polygamma factor.

\subsection{Compatibility of successive renormalized convolutions}
Theorem~\ref{thm:collision} proves compatibility of one collision with a
specified finite part. It does not establish associativity of successive
renormalized products. Determine whether an associative convolution
operation can be defined on a suitable space of periodic distributions,
and identify any finite local terms needed to recover the present
normalization.

\subsection{Arithmetic traces beyond the linear closure}
Character traces of the bilinear family reduce to ordinary Dirichlet
$L$-jets. Which additional coordinates survive for trilinear products,
unequal dilations, or nonrational shifts? Any claim of irreducibility
should be stated in an explicit formal law quotient unless numerical
period independence has actually been proved.

\subsection{Priority comparison, formalization and exact certificates}
Compare the all-order coefficient theorem systematically with the wider
Hurwitz-transform and finite-part-integral literature. A result already
implicit in a classical source should receive appropriate attribution,
while the explicit counterterm and coefficient presentation can still
serve as a useful repository contribution. Formalization can then
separate the local analytic lemmas, Fourier orthogonality, meromorphic
continuation and finite homogeneous algebra. The last layer already has
replayable exact data, but the present package is not a proof-assistant
certificate.

The existing colored-polylogarithm problems, including the next mixed
Gaussian reduction after $S_4$, remain distinct. They may benefit from
spectral generating functions, but no implication from this closure
theorem to their resolution is asserted.

\appendix
\section{Reproduction commands and artifact map}
\label{app:reproduce}

From the package root, the mathematical checks are run by
\begin{verbatim}
python code/exact_coefficients.py --max-total 8
python code/verify_numeric.py --dps 45
python code/verify_stieltjes_pairs.py
python code/polylog_derivative_canary.py
\end{verbatim}
The tested dependencies are SymPy 1.14.0 and mpmath 1.3.0. The scripts
write JSON result records and return a nonzero exit code when a required
identity check fails; the version-specific canary records its observation
without assuming that future software must reproduce the same loss.
For the PDF, use
\begin{verbatim}
latexmk -pdf -interaction=nonstopmode -halt-on-error article.tex
\end{verbatim}
The Makefile exposes these as \texttt{make check} and \texttt{make pdf}.
The integration fragment defines its own auxiliary macro names and uses
the label prefix \texttt{corrclosure:}. It is a proposed mathematical
section, not an automatic patch against a moving repository snapshot.
The article itself is self-contained and requires no external figure or
bibliography files.

\section{Coefficient data conventions}

Each row of \file{results/coefficients.json} contains the integer indices
\texttt{m}, \texttt{n}, a constant polynomial, an array \texttt{a} of
coefficients of $\gamma_j(a)$, an array \texttt{b} of coefficients of
$\gamma_j(1-a)$, and a polynomial \texttt{collision}.
Within these strings, \texttt{z2}, \texttt{z3}, and so on denote Riemann
zeta values; \texttt{L} denotes $\log a$ only in the collision polynomial.
This differs from the local notation $L=\log(2\pi)$ used in
Section~\ref{sec:loggamma}. The distinction is explicit in the JSON
schema and must be preserved when importing the data.

The symbolic table intentionally retains expressions such as
$\zeta(2)^2$ and $\zeta(4)$ separately. One may subsequently apply
$\zeta(2)^2=\tfrac52\zeta(4)$, but doing so changes the presentation
rather than the theorem. The coefficient generator does not use any
unproved independence conjecture to certify the rows.

\begin{thebibliography}{99}
\bibitem{dlmf-hurwitz}
NIST Digital Library of Mathematical Functions,
\emph{Section 25.11: Hurwitz Zeta Function}, particularly the shift,
Fourier, distribution, special-value and derivative identities.
\url{https://dlmf.nist.gov/25.11}. Accessed 10 October 2026.

\bibitem{dlmf-polylog}
NIST Digital Library of Mathematical Functions,
\emph{Section 25.12: Polylogarithms}.
\url{https://dlmf.nist.gov/25.12}. Accessed 10 October 2026.

\bibitem{shpot}
M.~A. Shpot, M.~P. Chaudhary and R.~B. Paris,
\emph{Integrals of products of Hurwitz zeta functions and the Casimir
effect in $\phi^4$ field theories}, arXiv:1603.00722, version 2.
\url{https://arxiv.org/html/1603.00722v2}.

\bibitem{repo-integration}
ProveIt contributors,
\emph{Zeta jets, Stieltjes antiderivatives and negative polygamma},
chapter source \file{07-integration.tex} in
\file{Analysis/Polylogarithms/docs/manuscript/chapters/}.
Selected sections inspected 10 October 2026.
\url{https://github.com/VladimirReshetnikov/ProveIt/tree/main/Analysis/Polylogarithms/docs/manuscript}.

\bibitem{repo-s4}
ProveIt contributors,
\emph{The mixed Gaussian weight-five identity and its exact proof},
\file{04-S4-proof.tex}, blob
\file{7ed3a5f375c027cc87a9c90d9f50e35ca447f1d2}, in the same manuscript.
Selected sections inspected 10 October 2026.

\bibitem{repo-discovery}
ProveIt contributors,
\emph{Experimental discovery and the route to proof},
\file{10-discovery.tex}, blob
\file{207c719ae6e7fb1639401f449216ae9b695afd40}, in the same manuscript.
Selected sections inspected 10 October 2026.

\bibitem{repo-incoming}
ProveIt contributors,
\emph{Incoming reports: intake instructions},
\file{docs/incoming/README.md}, blob
\file{9956ad83c5ddfa3c75e5f2946481a83082bc81e8}.
\url{https://github.com/VladimirReshetnikov/ProveIt/tree/main/docs/incoming}.
Selected intake sections and directory inventory inspected 10 October 2026.
\end{thebibliography}
\end{document}
